\chapter{Metodología y desarrollo}
\label{cap:metodologia}

Este capítulo describe cómo se ha realizado el trabajo. Debe permitir que otra
persona entienda el procedimiento seguido, los criterios de diseño, las
herramientas empleadas y las decisiones técnicas más importantes.

\section{Metodología empleada}
\label{sec:metodologia-empleada}

Explica el enfoque general del TFG: experimental, analítico, de diseño,
documental, computacional, de simulación, de implantación, comparativo o una
combinación de varios. La metodología debe estar conectada con los objetivos de
la Sección~\ref{sec:objetivos}.

Un esquema habitual para trabajos de ingeniería es el siguiente:

\begin{enumerate}
    \item Revisión bibliográfica, normativa y técnica.
    \item Definición de requisitos y alcance.
    \item Diseño de la solución o planteamiento del estudio.
    \item Desarrollo, cálculo, simulación, montaje o implementación.
    \item Pruebas, validación y análisis de resultados.
    \item Redacción de conclusiones y trabajo futuro.
\end{enumerate}

\subsection{Fases del estudio}
\label{subsec:fases-estudio}

La Tabla~\ref{tab:plan-trabajo} muestra un ejemplo de cómo presentar una
planificación por fases. Sustituye los textos de ejemplo por las fases reales
del proyecto.

\begin{table}[!ht]
    \footnotesize
    \centering
    \caption[Ejemplo de planificación del TFG por fases]{Ejemplo de planificación del TFG por fases. Fuente: elaboración propia}
    \label{tab:plan-trabajo}
    \begin{tabularx}{\textwidth}{p{3.4cm}X}
        \toprule
        \textbf{Fase} & \textbf{Descripción} \\
        \midrule
        Análisis inicial & Identificación del problema, revisión de documentación y delimitación del alcance. \\
        Diseño & Definición de requisitos, alternativas de solución y criterios de evaluación. \\
        Desarrollo & Implementación, cálculos, simulaciones, montaje o elaboración de la propuesta técnica. \\
        Validación & Pruebas, comparación de resultados, análisis de errores y discusión de limitaciones. \\
        Redacción & Integración de resultados, conclusiones, anexos y revisión formal de la memoria. \\
        \bottomrule
    \end{tabularx}
\end{table}

\subsection{Herramientas y materiales}
\label{subsec:herramientas-materiales}

Describe los recursos empleados: software, equipos de laboratorio, sensores,
instrumentación, bases de datos, lenguajes de programación, paquetes de cálculo,
normativa o documentación técnica. Indica versiones cuando sean relevantes para
la reproducibilidad.

La estructura recomendada para este apartado es:

\begin{itemize}
    \item Herramientas de diseño, cálculo o simulación.
    \item Materiales, equipos o componentes utilizados.
    \item Datos de partida y criterios de selección.
    \item Configuración del entorno de pruebas o desarrollo.
\end{itemize}

\section{Desarrollo de la propuesta}
\label{sec:desarrollo-propuesta}

En esta sección se explica el trabajo realizado paso a paso. Debe contener la
parte técnica central de la memoria: cálculos, diseños, diagramas, modelos,
procedimientos, implementación o montaje.

Si se incluyen ecuaciones, deben estar numeradas sólo cuando se vayan a
referenciar. Por ejemplo, la Ecuación~\ref{eq:variacion-relativa} muestra una
forma sencilla de calcular una variación relativa porcentual:

\begin{equation}
    \Delta_{\%} =
    \frac{V_{\mathrm{final}} - V_{\mathrm{inicial}}}{V_{\mathrm{inicial}}}
    \cdot 100
    \label{eq:variacion-relativa}
\end{equation}

Define cada variable antes o después de la ecuación y especifica sus unidades.
No dejes fórmulas aisladas sin interpretación técnica.

\subsection{Ejemplo de listado de código}
\label{subsec:ejemplo-listado}

Cuando el TFG incluya programación, configuración o pseudocódigo, puedes usar
el entorno \texttt{lstlisting}. El Listado~\ref{lst:ejemplo-validacion} muestra
un ejemplo mínimo.

\begin{lstlisting}[language=Python,caption={Ejemplo de validación de resultados},label={lst:ejemplo-validacion}]
def error_relativo(valor_medido, valor_referencia):
    """Calcula el error relativo porcentual."""
    return abs(valor_medido - valor_referencia) / valor_referencia * 100

resultado = error_relativo(9.8, 10.0)
print(f"Error relativo: {resultado:.2f} %")
\end{lstlisting}

No incluyas grandes bloques de código si no son necesarios para entender el
trabajo. Cuando el código sea extenso, es preferible explicar los fragmentos
relevantes y llevar el resto a un anexo o a un repositorio asociado.

\subsection{Ejemplo de algoritmo}
\label{subsec:ejemplo-algoritmo}

Los algoritmos ayudan a describir procedimientos de decisión, búsqueda,
optimización o validación sin depender de un lenguaje de programación concreto.

\begin{algorithm}[H]
    \DontPrintSemicolon
    \KwData{Requisitos del problema y alternativas disponibles}
    \KwResult{Alternativa seleccionada}
    Definir criterios de evaluación\;
    Asignar una ponderación a cada criterio\;
    \ForEach{alternativa}{
        Calcular la puntuación ponderada\;
    }
    Seleccionar la alternativa con mayor puntuación justificada\;
    \caption{Ejemplo de procedimiento de selección}
    \label{alg:seleccion-alternativa}
\end{algorithm}

\section{Consideraciones éticas, de seguridad y sostenibilidad}
\label{sec:consideraciones}

Incluye aquí cualquier consideración relevante sobre seguridad, protección de
datos, riesgos laborales, impacto ambiental, accesibilidad, sostenibilidad,
mantenimiento o uso responsable de los recursos.

Aunque el TFG sea principalmente técnico, es recomendable explicar las
restricciones no puramente técnicas que afectan a la solución. Este apartado
puede ser breve, pero debe demostrar que se han tenido en cuenta los efectos del
trabajo más allá del cálculo o del prototipo.

\section{Limitaciones metodológicas}
\label{sec:limitaciones-metodologicas}

Todo TFG tiene límites: tiempo, presupuesto, disponibilidad de equipos,
incertidumbre de datos, tamaño de muestra, alcance de las pruebas o supuestos de
cálculo. Explicarlos no debilita el trabajo; al contrario, ayuda a interpretar
correctamente los resultados.

Algunos ejemplos de limitaciones que puedes adaptar son:

\begin{enumerate}
    \item Disponibilidad limitada de datos experimentales o de documentación técnica.
    \item Restricciones de tiempo para realizar pruebas de larga duración.
    \item Uso de un entorno de laboratorio que no reproduce todas las condiciones reales.
    \item Dependencia de versiones concretas de software, equipos o librerías.
\end{enumerate}
